% =====================================================================
%  Replacement for \section{Transformers}.
% =====================================================================

% ---- local notation -------------------------------------------------
\providecommand{\Facc}{\F'}                                   % accumulator precision
\providecommand{\avghardmax}{\mathrm{avghardmax}}
\providecommand{\Wq}{W_{\textnormal{Q}}}
\providecommand{\Wk}{W_{\textnormal{K}}}
\providecommand{\Wv}{W_{\textnormal{V}}}
\providecommand{\Wout}{W_{\textnormal{out}}}
\providecommand{\bitof}[2]{\langle #1\rangle_{#2}}            % b-th binary digit
% counting quantifiers:  \cntge{k}{\phi}  =  "at least k positions y <= x satisfy phi"
\providecommand{\cntge}[2]{\exists^{\geq #1} y\leq x.\, #2}
\providecommand{\cntgestrict}[2]{\exists^{\geq #1} y< x.\, #2}
\providecommand{\cnteq}[2]{\exists^{= #1} y\leq x.\, #2}

\section{Transformers}
\label{sec:transformers}

We define the transformers we verify, and then translate each of them into a $\brasp$
autoregressor (\cref{sec:autoregressive-programs}) with the same next-symbol distributions.

\subsection{\texorpdfstring{$(1,1)$}{(1,1)}-precision transformers}
\label{sec:transformer_definitions}

Our transformers compute in fixed precision.
Let $\numbits\geq\fracbits+2$.
We write $\F$ for the set of numbers $m\cdot2^{-\fracbits}$ with $m\in\Z$ and
$-2^{\numbits-1}\leq m<2^{\numbits-1}$, that is, numbers with $\numbits$ bits, of which $\fracbits$ are fractional.
For $x\in\R$, $\round_{\F}(x)$ is the greatest element of $\F$ that is $\leq x$ (or $\min\F$ if there is none).
Rounding saturates: $\round_{\F}(x)=\max\F$ whenever $x\geq\max\F$.

We follow the $(1,1)$-precision transformers of \citet{li-cotterell-2025-characterizing,yang-etal-2025-length},
in which attention weights are rounded as well as activations.
The only change is that the normalizer of attention is accumulated in
$\Facc := \F_{2\numbits,\fracbits}$, which has $\numbits$ more integer bits than $\F$.
Rounding the normalizer into $\F$ itself would let unboundedly many weights equal $1$,
and the model could then compare counts.

\begin{remark}
Fixed precision is not an artificial restriction. Deployed transformers run in finite-precision formats
(e.g., 16-bit floating point or 8-bit integers), so every activation and attention weight is rounded to a
finite set of values, just as in $\F$. Accumulating the normalizer in a wider format than the activations is
also standard practice: matrix-multiplication kernels routinely sum in higher precision and round the result
back down. Our $\Facc$ models this.
\end{remark}

The following definition abstracts away from some details, but suffices for our purposes.

\begin{definition} \label{def:transformer}
A \defn{\frtfr} of depth $L$ and width $d$ consists of a word embedding $E\colon\alphabet\to\F^d$ and,
for each layer $\ell\in[L]$, functions $\Wq^{(\ell)},\Wk^{(\ell)},\Wv^{(\ell)},f^{(\ell)}\colon\F^d\to\F^d$.
On input $\str=w_1\cdots w_n$, it computes hidden states $\mathbf{h}^{(\ell)}_i\in\F^d$ for $i\in[n]$:
$\mathbf{h}^{(0)}_i = E(w_i)$, and for $\ell\in[L]$,
\begin{align}
A_{ij} &= \round_{\F}\bigl(\exp\bigl(\Wq^{(\ell)}(\mathbf{h}^{(\ell-1)}_i)\cdot\Wk^{(\ell)}(\mathbf{h}^{(\ell-1)}_j)\bigr)\bigr), &
B_i &= \round_{\Facc}\Bigl(\textstyle\sum_{j\leq i}A_{ij}\Bigr), \notag\\
\alpha_{ij} &= \round_{\F}(A_{ij}/B_i), &
\mathbf{c}_i &= \round_{\F}\Bigl(\textstyle\sum_{j\leq i}\alpha_{ij}\,\Wv^{(\ell)}(\mathbf{h}^{(\ell-1)}_j)\Bigr), \label{eq:attention}\\
\mathbf{h}^{(\ell)}_i &= f^{(\ell)}\bigl(\round_{\F}(\mathbf{c}_i+\mathbf{h}^{(\ell-1)}_i)\bigr), \notag
\end{align}
where $A_{ij},B_i,\alpha_{ij},\mathbf{c}_i$ depend on $\ell$, and $x/0:=0$.
\end{definition}

The position-wise functions are arbitrary, so they include affine maps, ReLU networks, and layer
normalization, all evaluated in fixed precision. Since attention ranges over $j\leq i$,
$\mathbf{h}^{(\ell)}_i$ depends only on $\str_{\leq i}$.

The property that makes these transformers translatable is that attention is spread over only
boundedly many positions.

\begin{lemma} \label{lem:bounded_support}
For every $i$, at most $2^{\fracbits}$ positions $j\leq i$ have $\alpha_{ij}\neq0$.
\end{lemma}

\begin{proof}
Every $A_{ij}$ is a nonnegative multiple of $2^{-\fracbits}$ below $2^{\numbits-1-\fracbits}$.
If the sum $\sum_{j\leq i}A_{ij}$ is less than $\max\Facc$, then it is represented exactly, so $B_i$ equals it and $\sum_j\alpha_{ij}\leq1$.
Each nonzero $\alpha_{ij}$ is at least $2^{-\fracbits}$, so at most $2^{\fracbits}$ of them are nonzero.
Otherwise $B_i=\max\Facc\geq2^{\numbits-1}$, so $A_{ij}/B_i<2^{-\fracbits}$ and every $\alpha_{ij}=0$.
\end{proof}

\subsection{Autoregressive transformers}
\label{sec:hardmax_autoregressive}

To generate strings, we add a readout that turns the last hidden state into a distribution over $\eos$.

\begin{definition} \label{def:hardmax_ar}
An \defn{autoregressive \frtfr} $\autoregressor=(\tfr,\Wout,\pi)$ consists of a \frtfr{} $\tfr$ of depth $L$,
a readout $\Wout\colon\F^d\to\F^{\eos}$, and a first-symbol distribution $\pi\in\Delta(\eos)$. For nonempty $\stru$ and $\tau\in\eos$,
\begin{equation}
\symprob{\autoregressor}{\stru}{\tau} := \avghardmax\bigl(\Wout(\mathbf{h}^{(L)}_{|\stru|})\bigr)(\tau),
\qquad
\avghardmax(\mathbf{z})(\tau) := \frac{\ind{\tau\in\argmax\mathbf{z}}}{|\argmax\mathbf{z}|}.
\label{eq:hardmax_symprob}
\end{equation}
\end{definition}

So each next-symbol distribution is uniform over the symbols with maximal logit.
As for $\brasp$ autoregressors (\cref{def:autoregressor}), $\symprob{\autoregressor}{\epsilon}{\tau}:=\pi(\tau)$,
and suffix and length-capped probabilities are defined as in \cref{eq:sufprob} and \cref{sec:verification-problem}.

\paragraph{Other readouts.}
The translation below uses only the fact that the next-symbol distribution is a function of the logits $\Wout(\mathbf{h}^{(L)}_{|\stru|})$.
So $\avghardmax$ can be replaced by any readout $\rho\colon\F^{\eos}\to\Delta(\eos)$.
The choice trades faithfulness to sampling against the cost of weighted model counting.
Each distinct transition probability needs its own weighted variable at every step (\cref{sec:bounded-brasp}),
so readouts with few distinct probabilities are cheaper. \Cref{tab:readouts} lists three options.
In the second row, $z_{(k)}$ is the $k$-th largest logit. In the third, each probability of $\softmax(\mathbf{z}/T)$
is rounded down to a multiple of $2^{-t}$, and the remaining mass goes to the first symbol in $\argmax\mathbf{z}$.

\begin{table}[h]
\caption{Readouts $\rho(\mathbf{z})$ and the number of distinct probabilities they produce.}
\label{tab:readouts}
\begin{tabular}{lll}
\toprule
Readout & Next-symbol distribution & Distinct probabilities \\
\midrule
$\avghardmax(\mathbf{z})$ & uniform over $\argmax\mathbf{z}$ & $1/m$ for $m\leq|\eos|$ \\
$\mathrm{top}_k(\mathbf{z})$ & uniform over $\{\tau\mid z_\tau\geq z_{(k)}\}$ & $1/m$ for $m\leq|\eos|$ \\
$\softmax_t(\mathbf{z}/T)$ & softmax at temperature $T$, $t$ bits & at most $2^t+1$ values \\
\bottomrule
\end{tabular}
\end{table}

\subsection{From transformers to \texorpdfstring{$\brasp$}{B-RASP}}
\label{sec:transformer_to_brasp}

We translate transformers into $\brasp$ in the style of
\citet[App.~B]{yang-etal-2025-knee} and \citet[App.~B.3]{yang-etal-2025-length}.
Those translations target logics with counting, so we first show how $\brasp$ can count
up to a constant, and how it can then compute sums over positions.

\paragraph{Counting.}
The \defn{counting quantifiers} are defined by induction on $k\geq 0$:
\begin{align*}
\cntgestrict{0}{\phi(y)} &:= \top &
\cntgestrict{k+1}{\phi(y)} &:= \attrdefault{y}{y<x}{\phi(y)}{\bigl(\exists^{\geq k} x<y.\,\phi(x)\bigr)}{\bot} \\
\cntge{0}{\phi(y)} &:= \top &
\cntge{k+1}{\phi(y)} &:= \cntgestrict{k+1}{\phi(y)} \lor \bigl(\phi(x)\land\cntgestrict{k}{\phi(y)}\bigr)
\end{align*}
and $\cnteq{k}{\phi(y)} := \cntge{k}{\phi(y)}\land\lnot\cntge{k+1}{\phi(y)}$.
By induction on $k$, $\str,i\models\cntge{k}{\phi(y)}$ iff at least $k$ positions $j\leq i$ satisfy $\phi$:
the attention finds the rightmost $j<i$ satisfying $\phi$, and the positions before $j$ must supply the other $k-1$.

\paragraph{Sums.}
Suppose that for a length-preserving $F\colon\alphabet^*\to(\F^d)^*$ we have formulas
$\phi_{F=\mathbf{u}}$, for $\mathbf{u}\in\F^d$, such that
\[
\str,i\models\phi_{F=\mathbf{u}} \iff F(\str)_i=\mathbf{u}.
\]
Then every function of $F$ is definable by a disjunction: for $g\colon\F^d\to X$ and $x\in X$,
\begin{equation}
\phi_{g(F)=x} := \bigvee_{\mathbf{u}\colon g(\mathbf{u})=x}\phi_{F=\mathbf{u}}.
\label{eq:pointwise}
\end{equation}
We will use this repeatedly. Sums over positions need more work. As in
\citet{yang-etal-2025-knee}, we sum a nonnegative integer quantity by counting, for each bit index,
the positions where that bit is $1$.
Writing $\bitof{n}{b}$ for the $b$-th binary digit of $n\in\N$,
\[
\textstyle\sum_{j\leq i} g(F(\str)_j) = \sum_b 2^b\cdot\bigl|\{j\leq i\mid \bitof{g(F(\str)_j)}{b}=1\}\bigr|.
\]
Each count only needs to be known up to a cap $K$, because any count of $K$ or more
already makes the sum at least $K$.

\begin{lemma} \label{lem:sums}
Let $g\colon\F^d\to\{0,\ldots,2^m-1\}$ and $K\geq1$. For each $k\leq K$ there is a formula
$\phi_{\Sigma g=k}$ such that
\[
\str,i\models\phi_{\Sigma g=k}\iff \min\Bigl(K,\textstyle\sum_{j\leq i}g(F(\str)_j)\Bigr)=k.
\]
\end{lemma}

\begin{proof}
Take a disjunction over the possible capped counts $n_b\in[0,K]$ for each bit index $b<m$:
\[
\phi_{\Sigma g=k} := \bigvee_{\substack{n_0,\ldots,n_{m-1}\in[0,K]\\ \min(K,\,\sum_b 2^b n_b)=k}}\ \bigwedge_{b<m}
\begin{cases}
\cnteq{n_b}{\phi_{\bitof{g(F)}{b}=1}(y)} & n_b<K\\
\cntge{K}{\phi_{\bitof{g(F)}{b}=1}(y)} & n_b=K.
\end{cases}
\]
Exactly one disjunct holds, namely the one in which each $n_b$ is the true count capped at $K$.
If no count is capped, the two sums are equal; otherwise both are at least $K$.
\end{proof}

We can now translate a transformer layer by layer.

\begin{theorem} \label{thm:transformer_to_brasp}
Let $\tfr$ be a \frtfr{} of depth $L$ and width $d$. For every $\ell\leq L$ and $\mathbf{u}\in\F^d$
there is a $\brasp$ formula $\phi_{\mathbf{h}^{(\ell)}=\mathbf{u}}$ such that
\begin{equation}
\str,i\models\phi_{\mathbf{h}^{(\ell)}=\mathbf{u}} \iff \mathbf{h}^{(\ell)}_i(\str)=\mathbf{u}.
\label{eq:simulates}
\end{equation}
These formulas have total size $|\alphabet|\cdot2^{\numbits d}+L\cdot2^{O(\numbits(\numbits+d))}$.
\end{theorem}

\begin{proof}
The construction proceeds by induction on $\ell$. For $\ell=0$, define
\[
\phi_{\mathbf{h}^{(0)}=\mathbf{u}} := \bigvee_{\sigma\colon E(\sigma)=\mathbf{u}}Q_\sigma.
\]
Now suppose that \cref{eq:simulates} holds for layer $\ell-1$, and consider layer $\ell$.
The summands of $B_i$ and $\mathbf{c}_i$ in \cref{eq:attention} depend on two positions ($i$ and $j$),
whereas a unary formula depends on one. But $\mathbf{h}^{(\ell-1)}_i$ can take only finitely many values,
and so can $B_i$, so we can enumerate them. Fix $\mathbf{q}\in\F^d$ and $\beta\in\Facc$, and consider
positions $i$ with $\mathbf{h}^{(\ell-1)}_i=\mathbf{q}$ and $B_i=\beta$.

The denominator $B_i$ is the sum of $\round_{\F}\bigl(\exp(\Wq(\mathbf{q})\cdot\Wk(\mathbf{h}^{(\ell-1)}_j))\bigr)$
over $j\leq i$, capped at $\max\Facc$.
Multiplying by $2^{\fracbits}$ makes the summands integers, so \cref{lem:sums} gives a formula
$\phi^{\mathbf{q}}_{B=\beta}$ that, for $\mathbf{h}^{(\ell-1)}_i=\mathbf{q}$, holds iff $B_i=\beta$.

Each coordinate $\kappa$ of the numerator is the rounding of
$\sum_{j\leq i}t(\mathbf{h}^{(\ell-1)}_j)$, where
\[
t(\mathbf{u}) := \round_{\F}\bigl(\round_{\F}(\exp(\Wq(\mathbf{q})\cdot\Wk(\mathbf{u})))/\beta\bigr)\cdot\Wv(\mathbf{u})_\kappa .
\]
Multiplying by $2^{2\fracbits}$ makes $t$ integer-valued. We sum its positive part $t^+$ and its negative part $t^-$ separately.
By \cref{lem:bounded_support}, at most $2^{\fracbits}$ summands are nonzero, so with
$K=2^{2\numbits+\fracbits}$ the cap of \cref{lem:sums} is never reached. Hence, for $\mathbf{z}\in\F^d$,
\[
\phi^{\mathbf{q},\beta}_{\mathbf{c}=\mathbf{z}} := \bigwedge_{\kappa\in[d]}\ \bigvee_{\substack{k^+,k^-\colon\\ \round_{\F}(2^{-2\fracbits}(k^+-k^-))=z_\kappa}}
\phi_{\Sigma 2^{2\fracbits}t^+=k^+}\land\phi_{\Sigma 2^{2\fracbits}t^-=k^-}
\]
holds iff $\mathbf{c}_i=\mathbf{z}$.
Finally, taking the disjunction over all $\mathbf{q}$ and $\beta$,
\[
\phi_{\mathbf{h}^{(\ell)}=\mathbf{u}} := \bigvee_{\substack{\mathbf{q},\mathbf{z},\beta\colon\\ f^{(\ell)}(\round_{\F}(\mathbf{z}+\mathbf{q}))=\mathbf{u}}}
\phi_{\mathbf{h}^{(\ell-1)}=\mathbf{q}}\land\phi^{\mathbf{q}}_{B=\beta}\land\phi^{\mathbf{q},\beta}_{\mathbf{c}=\mathbf{z}}.
\]
At each position, exactly one choice of $\mathbf{q}$ and $\beta$ satisfies the first two conjuncts, namely the true values.
So \cref{eq:simulates} holds for layer $\ell$.

For the size, each use of \cref{lem:sums} has $m=O(\numbits)$ and $K=2^{O(\numbits)}$, and so has size
$2^{O(\numbits^2)}\cdot2^{\numbits d}$. There are $2^{O(\numbits d)}$ choices of $\mathbf{q},\beta,\mathbf{z},\kappa$.
\end{proof}

The readout is a function of $\mathbf{h}^{(L)}$, so by \cref{eq:pointwise} it becomes one formula per output symbol.

\begin{corollary} \label{cor:hardmax_to_brasp_ar}
Let $\autoregressor=(\tfr,\Wout,\pi)$ be an autoregressive \frtfr{}, and let
$\autoregressor'=((\phi_\tau)_{\tau\in\eos},\arout,\pi)$, where
\[
\phi_\tau := \phi_{\ind{\tau\in\argmax\Wout(\mathbf{h}^{(L)})}=1},
\qquad
\arout(\mathbf{b})(\tau) := \frac{\ind{\mathbf{b}_\tau=\top}}{|\{\tau'\mid\mathbf{b}_{\tau'}=\top\}|}.
\]
Then $\symprob{\autoregressor'}{\stru}{\tau}=\symprob{\autoregressor}{\stru}{\tau}$ for all $\stru\in\alphabet^+$ and $\tau\in\eos$.
So $\autoregressor$ and $\autoregressor'$ have the same suffix and length-capped probabilities, and therefore
the same answers to $\mathsf{PROB}$ and $\mathsf{THRESH}$.
\end{corollary}

\begin{proof}
$\phi_\tau$ holds at $|\stru|$ iff $\tau\in\argmax\Wout(\mathbf{h}^{(L)}_{|\stru|})$, and $\arout$ is uniform over exactly these symbols.
\end{proof}

The same construction works for $\mathrm{top}_k$, with $\phi_\tau := \phi_{\ind{\tau\in\mathrm{supp}(\mathrm{top}_k(\Wout(\mathbf{h}^{(L)})))}=1}$.
For an arbitrary readout $\rho$, we use one formula per next-symbol distribution instead of one per symbol.
Let $R:=\{\rho(\Wout(\mathbf{u}))\mid\mathbf{u}\in\F^d\}$ be the finite set of distributions that the readout can produce, and take
\[
\phi_\mu := \phi_{\rho(\Wout(\mathbf{h}^{(L)}))=\mu}\quad(\mu\in R),
\qquad
\arout(\mathbf{b}) := \text{the } \mu \text{ with } \mathbf{b}_\mu=\top.
\]
Exactly one $\phi_\mu$ holds at each position, so $\arout$ is well defined, and the proof is the same.
The transition probabilities of the resulting autoregressor are the values taken by distributions in $R$.